\section{Mechanization}
\label{sec:mechanization}

The decision procedures and their correctness are mechanized in Lean~4
\cite{MouraUllrich21} with Mathlib \cite{Mathlib20}. The development defines
trees over $\BFT$ and over $\BF$, the order with variance, the three classes
of trees, terms, quantifier-free formulas, $\Sigma_1$ sentences and their
Boolean combinations, and it contains
\begin{itemize}
\item a decision procedure for Boolean combinations of $\Sigma_1$ and $\Pi_1$
  sentences over each class, for every arity and variance over $\BFT$ and
  for a covariant constructor over $\BF$, with the equivalence of truth over
  regular and over arbitrary trees (\cref{thm:main});
\item decision procedures for disjunctive entailment between flat
  constraints over the three classes and both signatures, and the statements
  that satisfiability and (disjunctive) entailment are the sentences of
  \cref{sec:prelim-constraints};
\item the small-hole theorem for structural readers over an arbitrary finite
  alphabet (\cref{thm:small-hole}), with the computable bound $\Hole$, and
  its closed form (\cref{thm:elementary}).
\end{itemize}
The results of
\cref{sec:sentences,sec:readers,sec:overview,sec:fresh,sec:tails,sec:decision}
and the inequalities of \cref{thm:elementary} have counterparts in the
development; \cref{app:lean} gives the correspondence and the points where
the formal statements or proofs are arranged differently. The
complexity statements of \cref{sec:complexity} other than the numerical
bound, that is, \cref{cor:complexity}, \cref{prop:hardness} and
\cref{ex:lcm}, are not mechanized.

Each decision procedure is a computable definition together with a theorem
that it returns \emph{true} exactly when its proposition holds. The
procedures have no hypotheses beyond a finite type of variables with
decidable equality: satisfiability is tested inside them, the number of
queries may be zero, and the arity may be zero. Decidability comes from
bounded searches. Classical reasoning is used in the proofs of correctness
and of the bounds, where it is harmless, and not to choose answers. The
procedures are executable in principle and useless in practice, because of
the size of the bound; they establish decidability.

The statements one has to trust are the definitions of trees and of the
order, of the syntax and truth of sentences, of the classes of trees, of
disjunctive entailment, and the statements of the main theorems, which are
collected in one file. The kernel checks everything else. The
development uses no axioms beyond the three standard ones of Lean and
Mathlib (propositional extensionality, quotient soundness and choice), no
unfinished proofs, and no native evaluation; a file of the development
checks the axioms of every public declaration and that every decision
procedure compiles to executable code. Examples show that the definitions
are not vacuous: they include the entailments of \cref{ex:unary} and
\eqref{eq:arrow-disj} together with the failure of their single disjuncts,
sentences that separate finite from regular trees, a reader compiled from an
entailment with two queries, with a hole, and a minimal fresh tuple.

The chain of shortening arguments is long, and several of its steps are of a
kind in which informal proofs are fragile: the transport of comparisons
through deletions has many boundary cases, with cuts at the ends of deleted
factors and empty rests. Both deletion arguments are therefore built on one
small library for deleting a set of positions from a word and following cuts
and runs through the deletion (\cref{app:fresh-deletion,app:compression}).

The development has about $27{,}500$ lines. It is available at
\begin{center}
\url{https://github.com/noti0na1/deci-sigma1}.
\end{center}
